% =====================================================================
%  Chapter 7: Chiral symmetry on the lattice
%  Gattringer & Lang, "Quantum Chromodynamics on the Lattice" (LNP 788)
%  Reconstructed from extract/ch07.txt; equations auto-numbered.
% =====================================================================

\chapter{Chiral symmetry on the lattice}

Chiral symmetry and its spontaneous breaking are central properties of QCD
with important phenomenological implications. They explain why pions have
unexpectedly small masses and why we do not see degenerate masses for parity
partners in the baryon sector. This important role of chiral symmetry and the
mechanism of its breaking require that a reasonable lattice version of QCD
has to implement chiral symmetry properly.

However, putting chiral symmetry onto the lattice turned out to be a
formidable challenge. The underlying reason is the doubling problem. In
Chap.~5 we have added the Wilson term to the naive fermion action in order to
remove the doublers. This term, however, breaks chiral symmetry explicitly.
Only more than 20 years after Wilson's introduction of lattice gauge theory
the problem of chiral symmetry on the lattice was solved with a generalization
of chiral symmetry through the so-called ``Ginsparg--Wilson equation'' for the
lattice Dirac operator. With this new concept a clean implementation of chiral
symmetry on the lattice has been achieved.

This chapter is devoted to chiral symmetry on the lattice. After a brief
review of chiral symmetry in continuum QCD we discuss why Wilson fermions
fail to show chiral symmetry. Subsequently we present the Ginsparg--Wilson
equation and explore its consequences concerning the axial anomaly and the
index theorem. Finally we discuss an explicit solution of the Ginsparg--Wilson
equation, the so-called overlap operator.

\section{Chiral symmetry in continuum QCD}

\subsection{Chiral symmetry for a single flavor}

To warm up, let us first discuss chiral symmetry for the case of only a single
flavor in the continuum. The action for a massless fermion reads (for
notational convenience we drop the space--time argument $x$ in the following)
\begin{equation}
  S_F[\psi,\bar\psi,A] = \int d^4x\;\mathcal{L}(\psi,\bar\psi,A),
  \qquad
  \mathcal{L}(\psi,\bar\psi,A) = \bar\psi\,\gamma_\mu(\partial_\mu + iA_\mu)\,\psi
  = \bar\psi D\psi,
\end{equation}
where $D$ denotes the massless Dirac operator.

We now perform a chiral rotation of the fermion fields
\begin{equation}
  \psi \to \psi' = e^{i\alpha\gamma_5}\psi,
  \qquad
  \bar\psi \to \bar\psi' = \bar\psi\,e^{i\alpha\gamma_5},
\end{equation}
where $\gamma_5$ is the chirality matrix acting in Dirac space and $\alpha$ is
a constant, real parameter. The Lagrangian density is invariant under the
chiral rotation,
\begin{align}
  \mathcal{L}(\psi',\bar\psi',A)
  &= \bar\psi'\,\gamma_\mu(\partial_\mu + iA_\mu)\,\psi'
  = \bar\psi\,e^{i\alpha\gamma_5}\,\gamma_\mu(\partial_\mu + iA_\mu)\,e^{i\alpha\gamma_5}\,\psi \notag\\
  &= \bar\psi\,e^{i\alpha\gamma_5}\,e^{-i\alpha\gamma_5}\,\gamma_\mu(\partial_\mu + iA_\mu)\,\psi
  = \mathcal{L}(\psi,\bar\psi,A),
\end{align}
where in the second step we have used $\gamma_\mu\gamma_5 = -\gamma_5\gamma_\mu$.
A mass term breaks this invariance since it transforms nontrivially:
\begin{equation}
  m\,\bar\psi'\psi' = m\,\bar\psi\,e^{i2\alpha\gamma_5}\psi.
\end{equation}

The chiral symmetry decouples the action of left- and right-handed massless
fermions. To see this we introduce the right- and left-handed projectors
\begin{equation}
  P_R = \frac{1+\gamma_5}{2},
  \qquad
  P_L = \frac{1-\gamma_5}{2},
\end{equation}
which obey
\begin{align}
  P_R^2 = P_R,\qquad P_L^2 = P_L, \qquad P_R P_L = P_L P_R = 0, \qquad P_R + P_L = 1,
  \notag\\
  \gamma_\mu P_L = P_R\gamma_\mu, \qquad \gamma_\mu P_R = P_L\gamma_\mu .
\end{align}
With these we can define right- and left-handed fermion fields
\begin{equation}
  \psi_R = P_R\psi, \qquad \psi_L = P_L\psi,
  \qquad
  \bar\psi_R = \bar\psi P_L, \qquad \bar\psi_L = \bar\psi P_R .
\end{equation}
After a few lines of algebra the decoupling of left- and right-handed
components follows
\begin{equation}
  \mathcal{L}(\psi,\bar\psi,A) = \bar\psi_L D\psi_L + \bar\psi_R D\psi_R .
\end{equation}
Thus we find that left- and right-handed components ``do not talk to each
other.'' A mass term, however, mixes the components:
\begin{equation}
  m\,\bar\psi\psi = m\left(\bar\psi_R\psi_L + \bar\psi_L\psi_R\right).
\end{equation}
Since chiral symmetry of the action holds only for massless quarks, the limit
of vanishing quark mass is often referred to as the chiral limit.

One can summarize the essence of chiral symmetry in the simple equation
\begin{equation}
  D\gamma_5 + \gamma_5 D = 0 .
\end{equation}
It expresses the fact that the massless Dirac operator
$D = \gamma_\mu(\partial_\mu + iA_\mu)$ anticommutes with $\gamma_5$.

\subsection{Several flavors}

Let us now delve somewhat deeper into the symmetry properties of the QCD
action and consider the theory with $\Nf$ flavors of quarks. Now the fermion
fields $\psi$, $\bar\psi$ carry also a flavor index. However, for notational
convenience we use vector notation for the flavor index, too. In this notation
the fermion action reads
\begin{equation}
  S_F[\psi,\bar\psi,A] = \int d^4x\;\bar\psi\left(\gamma_\mu(\partial_\mu + iA_\mu) + M\right)\psi,
\end{equation}
where we introduce a mass matrix
\begin{equation}
  M = \mathrm{diag}(m_1, m_2, \ldots, m_{\Nf})
\end{equation}
acting in flavor space (in the real world $m_1 = m_u$, $m_2 = m_d$, etc.).
For generalizing the chiral transformations of the last paragraph we use
matrices that mix the different flavors. In particular we denote the
generators of $SU(\Nf)$ by $T_i$, $i = 1, 2, \ldots, \Nf^2 - 1$ (see
Appendix~A.1).

Before we discuss the chiral rotations, let us first note that the action for
massless quarks is also invariant under the $\Nf^2$ vector
transformations\footnote{The name ``vector transformations'' comes from the
fact that the corresponding Noether currents are vector currents.}
\begin{equation}
  \psi' = e^{i\alpha T_i}\psi,
  \qquad
  \bar\psi' = \bar\psi\,e^{-i\alpha T_i},
\end{equation}
\begin{equation}
  \psi' = e^{i\alpha\,\mathbb{1}}\psi,
  \qquad
  \bar\psi' = \bar\psi\,e^{-i\alpha\,\mathbb{1}},
\end{equation}
where $\mathbb{1}$ denotes the $\Nf \times \Nf$ unit matrix. The invariance of
(7.11) for $M = 0$ under transformations (7.13) and (7.14) is evident.
However, the invariance under (7.13) extends also to the case of degenerate
masses $M = \mathrm{diag}(m, m, \ldots, m)$. This symmetry is the isospin
symmetry, generalized to $\Nf$ flavors. Symmetry (7.14) even holds for
arbitrary masses and the corresponding conserved quantity is baryon number.

The chiral or axial vector rotations are defined as
\begin{equation}
  \psi' = e^{i\alpha\gamma_5 T_i}\psi,
  \qquad
  \bar\psi' = \bar\psi\,e^{i\alpha\gamma_5 T_i},
\end{equation}
\begin{equation}
  \psi' = e^{i\alpha\gamma_5\,\mathbb{1}}\psi,
  \qquad
  \bar\psi' = \bar\psi\,e^{i\alpha\gamma_5\,\mathbb{1}},
\end{equation}
where now also the left- and right-handed components of the different flavors
mix. Like for the case of a single flavor one finds that the action (7.11) is
invariant under (7.15) and (7.16) for the case of $M = 0$.

Taken altogether, the massless action has the symmetry
\begin{equation}
  SU(\Nf)_L \times SU(\Nf)_R \times U(1)_V \times U(1)_A .
\end{equation}
The notation used for the two $SU(\Nf)$ pieces stresses that the terms for the
left- and right-handed components are symmetric under independent $SU(\Nf)$
rotations.

However, when one considers the fully quantized theory one finds that the
fermion determinant is not invariant under (7.16) and the corresponding
so-called $U(1)_A$ axial symmetry is broken explicitly by a noninvariance of
the fermion integration measure (we will show this in Sect.~7.3). Taking into
account this so-called axial anomaly [1, 2] we find that for the quantized,
massless theory the symmetry is broken explicitly to the remaining symmetry
\begin{equation}
  SU(\Nf)_L \times SU(\Nf)_R \times U(1)_V .
\end{equation}
Introducing nonvanishing, degenerate masses, $M = \mathrm{diag}(m, m \ldots
m)$, breaks the symmetry $SU(\Nf)_L \times SU(\Nf)_R$ explicitly to its
subgroup $SU(\Nf)_V$,
\begin{equation}
  SU(\Nf)_V \times U(1)_V,
\end{equation}
and allowing for nondegenerate masses reduces the symmetry further to
\begin{equation}
  U(1)_V \times U(1)_V \times \ldots \times U(1)_V
  \qquad (\Nf \text{ factors}) .
\end{equation}

It is interesting to have a second look at the explicit breaking due to the
quark masses. The lightest quarks have masses that are relatively small
compared to the typical QCD scale of 1\,GeV given by, e.g., the mass of the
proton ($\approx$\,940\,MeV). In particular one has (in the $\overline{MS}$
renormalization scheme)
\begin{equation}
  m_u \approx m_d \approx 5\,\mathrm{MeV}, \qquad m_s \approx 100\,\mathrm{MeV},
\end{equation}
such that the explicit breaking of the chiral symmetry is very small for $u$
and $d$ quarks and a 10\% effect for the $s$ quark. Thus we would expect that
(7.18) is a good approximate symmetry for $\Nf = 2$ and in part also for
$\Nf = 3$.

\subsection{Spontaneous breaking of chiral symmetry}

If $u$ and $d$ quarks were massless, $SU(2)_L \times SU(2)_R \times U(1)_V$
would be an exact symmetry of QCD (for $\Nf = 2$) at least on the Lagrangian
level. Due to the small mass of the $u$ and $d$ quarks the explicit breaking
is very small and traces of the symmetry should be visible. One of the
consequences is that one expects degenerate masses for the nucleon and its
partner with negative parity, the so-called $N^*$.

We briefly sketch the argument for such a degeneracy of the masses: In the
last chapter we have shown that the masses of the nucleon and its negative
parity partner are obtained from the exponential decay of the zero momentum
correlators
\begin{equation}
  C_\pm(t) = -\int d^3x\,
  \left\langle
    \varepsilon_{abc}\,u(0)^T_a C\gamma_5 d(0)_b u(0)_c\,
    \tfrac{1\pm\gamma_4}{2}\,
    \varepsilon_{efg}\,\bar{u}_e(x) u(x)^T_f C\gamma_5 d(x)_g
  \right\rangle,
\end{equation}
where $x = (\vec{x}, t)$. This equation gives the continuum version of the
lattice correlator (6.21) after projection to zero momentum. The exponential
decay of $C_+(t)$ gives the mass of the nucleon (for forward running $t$) and
$C_-(t)$ corresponds to the $N^*$. Now consider the transformation
\begin{equation}
  \psi \to \exp\left(i\alpha\gamma_5 T_3\right)\psi,
  \qquad
  \bar\psi \to \bar\psi\exp\left(i\alpha\gamma_5 T_3\right),
\end{equation}
where $\psi$ and $\bar\psi$ are vectors in flavor space:
$\psi = (u, d)^T$, $\bar\psi = (\bar{u}, \bar{d})$. Using
$T_3 = \sigma_3/2 = \mathrm{diag}(1/2, -1/2)$ we find that the individual
flavors transform as
\begin{align}
  u &\to \exp\left(+i\,\frac{\alpha}{2}\gamma_5\right)u, \qquad
  \bar{u} \to \bar{u}\exp\left(+i\,\frac{\alpha}{2}\gamma_5\right), \notag\\
  d &\to \exp\left(-i\,\frac{\alpha}{2}\gamma_5\right)d, \qquad
  \bar{d} \to \bar{d}\exp\left(-i\,\frac{\alpha}{2}\gamma_5\right).
\end{align}
When setting $\alpha = \pi$ this simplifies to
\begin{equation}
  u \to i\gamma_5 u, \qquad
  \bar{u} \to i\bar{u}\gamma_5, \qquad
  d \to -i\gamma_5 d, \qquad
  \bar{d} \to -i\bar{d}\gamma_5 .
\end{equation}
Inserting this transformation into (7.22) one finds that the correlators
transform under (7.25) as
\begin{equation}
  C_\pm(t) \to -C_\mp(t) .
\end{equation}
Thus, up to an overall minus sign, the correlator $C_+(t)$ for the nucleon and
the correlator $C_-(t)$ for the $N^*$ are transformed into each other. Since
an overall minus sign does not alter the exponential decay properties of the
correlator, the invariance under (7.23) implies that the nucleon and its
parity partner $N^*$ have equal masses.

However, no such degeneracy of the nucleon and $N^*$ masses is observed. The
nucleon (proton or neutron) has a mass of 940\,MeV, while the $N^*$ is found
at 1535\,MeV. The mass difference of almost 600\,MeV is far too large to be
explained by the small explicit breaking of chiral symmetry due to the $u$ and
$d$ quark masses. So another mechanism which breaks the symmetry under (7.23)
must be at work. Since the action itself is invariant (up to the small
explicit breaking by the quark masses), the strong breaking effect observed in
nature must come from a spontaneous breaking of chiral symmetry, except for
$U(1)_A$ which is broken explicitly by the anomaly.

The spontaneous breaking of symmetries is a concept which is well known from,
e.g., spin systems describing ferromagnets. While the Hamiltonian of the
system is invariant under a global rotation of the spins, in the ground state
of the system all spins point in the same direction, thus leading to a
macroscopic magnetization of the system.

A similar effect is taking place in QCD. While the action is invariant under
chiral rotations, the ground state is not. An order parameter for chiral
symmetry breaking is
\begin{equation}
  \ev{u(x)\bar{u}(x)},
\end{equation}
the so-called chiral condensate. The chiral condensate transforms like a mass
term and is not invariant under chiral rotations. Consequently, when the
theory has a nonvanishing chiral condensate this implies that chiral symmetry
is broken spontaneously. In Sect.~7.3 we will discuss the calculation of the
chiral condensate on the lattice and show that it has a nonvanishing value.

Another important consequence of a spontaneously broken continuous symmetry
is the appearance of so-called Goldstone modes or Goldstone bosons (see
standard text books on quantum field theory such as [3, 4]). Goldstone modes
are massless bosonic excitations. For the example of the spin system with
continuous symmetry group these are spin waves. In the case of QCD the pions
are interpreted as the ``would-be'' Goldstone bosons of chiral symmetry
breaking. For massless quarks a spontaneous symmetry breaking mechanism of QCD
would explain massless pions. Their relatively small masses of about 140\,MeV
can be understood as resulting from the explicit breaking by the $u$ and $d$
quark masses.

To summarize, we find that for the light quark flavors the action is
essentially invariant under chiral transformations -- up to the chiral $U(1)$
transformation symmetry, which is broken explicitly by the axial anomaly.
However, the experimental findings do not reflect the chiral symmetry and we
must conclude that chiral symmetry is broken spontaneously due to the dynamics
of QCD. Chiral symmetry at the classical level and its spontaneous breaking
provide an explanation for several phenomena, such as the nondegenerate masses
of baryonic parity partners or the small masses of the pions.

\section{Chiral symmetry and the lattice}

Having convinced ourselves of the crucial importance of chiral symmetry and
its breaking for QCD phenomenology, let us now return to the lattice. We start
with analyzing the chiral properties of Wilson fermions.

\subsection{Wilson fermions and the Nielsen--Ninomiya theorem}

Upon inspecting the Wilson Dirac operator (5.51) we find that (7.10) is
violated also for vanishing quark mass. The culprit is the Wilson term (5.50)
which we had to add to the naive action in order to remove the doublers. This
term comes with the unit matrix $\mathbb{1}$ in Dirac space which does not
anticommute with $\gamma_5$ and therefore violates (7.10). Thus, even for
massless quarks we break chiral symmetry explicitly. Consequently we cannot
expect to capture such a subtle effect as spontaneous chiral symmetry breaking
with Wilson fermions.

This situation cannot be simply overcome by adding a term for removing the
doublers in a different way. There is a fundamental theorem by Nielsen and
Ninomiya [5--7] which states that on the lattice one cannot implement chiral
symmetry in the form of (7.10) and at the same time have a theory free of
doublers. More explicitly [8] consider a Euclidean action for free fermions in
the form
\begin{equation}
  S = \sum_{n,m,\mu} \bar\psi(n)\,i\gamma_\mu\,F_\mu(n|m)\,P_R\,\psi(m) .
\end{equation}
The kernel $F(n|m)$ is assumed to be translationally invariant,
$F_\mu(n|m) = F_\mu(n-m)$, and should give rise to a hermitian Hamilton
operator which implies $F_\mu(-n) = F_\mu(n)^*$. Furthermore, $F_\mu(n)$ must
be local, i.e., decrease sufficiently fast, such that its Fourier transform
$\hat{F}_\mu(p)$ exists and all its derivatives are continuous. Due to the
right-handed projector $P_R$ one would naively expect that action (7.28)
describes only right-handed fermions. However, the arguments in [8], based on
the expansion of $\hat{F}_\mu(p)$ around its zeros (poles in the propagator),
show that an equal number of left- and right-handed fermions is described by
an action of the form (7.28).

For many years the Nielsen--Ninomiya theorem seemed to define the ultimate
limitation for a further development of lattice QCD. A deep understanding of
chiral symmetry on the lattice became possible only after the rediscovery of
the Ginsparg--Wilson equation.

\subsection{The Ginsparg--Wilson equation}

In their seminal paper [9], Ginsparg and Wilson formulated the essential
equation for chiral symmetry on the lattice. Based on renormalization group
transformations (which we will discuss in Chap.~9) they proposed to replace
the continuum expression (7.10) by
\begin{equation}
  D\gamma_5 + \gamma_5 D = a\,D\gamma_5 D .
\end{equation}
In this equation the anti-commutator (7.10) is augmented with a nonvanishing
right-hand side. For dimensional reasons -- the Dirac operator $D$ has
dimension 1/length -- an extra factor of the lattice constant $a$ appears for
the term quadratic in $D$. Thus the right-hand side vanishes for $a \to 0$ and
the continuum form of chiral symmetry is recovered in a naive continuum limit.
However, (7.29) allows to define chiral symmetry on the lattice also for
finite $a$. We will discuss this construction and the consequences of the
Ginsparg--Wilson equation in great detail below.

It is interesting to analyze the effect of the extra term for the quark
propagator $D^{-1}$. Assuming that $D$ has no zero modes and is invertible, we
may multiply (7.29) with $D^{-1}$ from both sides and obtain
\begin{equation}
  \gamma_5 D^{-1}(n|m) + D^{-1}(n|m)\,\gamma_5 = a\,\gamma_5\,\delta(n-m),
\end{equation}
where we write explicitly the space--time arguments $n$, $m$ of the
propagators. The term introduced by Ginsparg and Wilson is simply a contact
term, i.e., the anti-commutator of the propagator with $\gamma_5$ is modified
only for $n = m$.

Before we explore the consequences of the Ginsparg--Wilson equation, we should
remark that the original paper [9] from 1982 had no immediate consequences and
was ``forgotten,'' since no solution of the nonlinear equation (7.29) was
imagined for the interacting case (nontrivial gauge links) at that time. Only
in 1997 and 1998 it was realized [10, 11] that two independent approaches to
chiral symmetry on the lattice, overlap and fixed point fermions introduced
earlier, give rise to Dirac operators $D$ obeying (7.29). We discuss these in
Sect.~7.4 and Chap.~9.

\subsection{Chiral symmetry on the lattice}

We have already announced that, based on the Ginsparg--Wilson equation, it is
possible to define a modified chiral rotation which leaves the lattice action
for massless quarks invariant [12]. Let $D$ be a lattice Dirac operator
obeying the Ginsparg--Wilson equation (7.29). Using $D$ we can define a chiral
rotation, which for $a \to 0$ reduces to the continuum transformation (7.2),
\begin{equation}
  \psi' = \exp\!\left(i\alpha\gamma_5\left(1-\frac{aD}{2}\right)\right)\psi,
  \qquad
  \bar\psi' = \bar\psi\exp\!\left(i\alpha\left(1-\frac{aD}{2}\right)\gamma_5\right).
\end{equation}
In this equation the action of the Dirac operator is understood in the sense
\begin{equation}
  (D\psi)(n)^a_\alpha = \sum_{m,\beta,b} D(n|m)^{ab}_{\alpha\beta}\,\psi(m)^b_\beta,
  \qquad
  (\bar\psi D)(m)^b_\beta = \sum_{n,\alpha,a} \bar\psi(n)^a_\alpha\,D(n|m)^{ab}_{\alpha\beta}.
\end{equation}
In the subsequent equations we refer to this convention, but use
matrix/vector notation for all indices: space--time, Dirac, and color.

The Lagrangian density for massless fermions for such a Dirac operator is
invariant under the transformation (7.31):
\begin{align}
  \mathcal{L}(\psi',\bar\psi') &= \bar\psi'\,D\,\psi' \notag\\
  &= \bar\psi\exp\!\left(i\alpha\left(1-\frac{aD}{2}\right)\gamma_5\right)D
     \exp\!\left(i\alpha\gamma_5\left(1-\frac{aD}{2}\right)\right)\psi \notag\\
  &= \bar\psi\exp\!\left(i\alpha\left(1-\frac{aD}{2}\right)\gamma_5\right)
     \exp\!\left(-i\alpha\left(1-\frac{aD}{2}\right)\gamma_5\right)D\,\psi \notag\\
  &= \bar\psi\,D\,\psi = \mathcal{L}(\psi,\bar\psi) .
\end{align}
In the step from the second to the third line we have used
\begin{equation}
  D\,\gamma_5\left(1-\frac{aD}{2}\right) + \left(1-\frac{aD}{2}\right)\gamma_5\,D = 0,
\end{equation}
which is just another way of writing the Ginsparg--Wilson equation (7.29).

Decomposition (7.7) into left- and right-handed components can be generalized
to the lattice. We define a new kind of projectors
\begin{equation}
  \hat P_R = \frac{1+\hat\gamma_5}{2}, \qquad
  \hat P_L = \frac{1-\hat\gamma_5}{2}, \qquad
  \hat\gamma_5 = \gamma_5(1 - aD) .
\end{equation}
Due to the Ginsparg--Wilson equation (7.29) one finds
$\hat\gamma_5^2 = 1$, implying
\begin{equation}
  \hat P_R^2 = \hat P_R, \qquad
  \hat P_L^2 = \hat P_L, \qquad
  \hat P_R\hat P_L = \hat P_L\hat P_R = 0, \qquad
  \hat P_R + \hat P_L = 1 .
\end{equation}
Invoking again the Ginsparg--Wilson equation one shows
\begin{equation}
  D\hat P_R = P_L D, \qquad
  D\hat P_L = P_R D .
\end{equation}
Based on this last equation a sensible definition of left- and right-handed
components, which generalizes the continuum form (7.7), is
\begin{equation}
  \psi_R = \hat P_R\psi, \qquad \psi_L = \hat P_L\psi,
  \qquad
  \bar\psi_R = \bar\psi P_L, \qquad \bar\psi_L = \bar\psi P_R .
\end{equation}
This definition, together with (7.37), implies the vanishing of terms like
$\bar\psi_L D\psi_R$ and $\bar\psi_R D\psi_L$. The action thus can be
decomposed like the continuum action (7.8) into left- and right-handed parts:
\begin{equation}
  \bar\psi D\psi = \bar\psi_L D\psi_L + \bar\psi_R D\psi_R .
\end{equation}
We thus may identify the continuum fields $\psi_L$, $\psi_R$ with the lattice
quantities (7.38). The symmetry breaking mass term (7.9) is identified as the
lattice term
\begin{equation}
  m\left(\bar\psi_R\psi_L + \bar\psi_L\psi_R\right)
  = m\,\bar\psi\left(P_L\hat P_L + P_R\hat P_R\right)\psi
  = m\,\bar\psi\left(1-\frac{aD}{2}\right)\psi .
\end{equation}
Thus Ginsparg--Wilson fermions with mass are described by the operator [13]
\begin{equation}
  D_m = D + m\left(1-\frac{aD}{2}\right) = \omega D + m\,\mathbb{1},
\end{equation}
with the abbreviation
\begin{equation}
  \omega \equiv 1-\frac{am}{2} .
\end{equation}
With chiral rotation (7.31) and decomposition (7.39) we have successfully
implemented the corresponding continuum structures (7.2) and (7.8) on the
lattice. The generalization to several flavors proceeds along the same steps
as for the continuum case.

Let us stress the most important difference between the concept of chirality
in the continuum and on the lattice. In the continuum chirality is a strictly
local concept independent of the gauge field. Rotation (7.2), as well as
projection (7.7), only involves the spinors at a given space--time point $x$,
i.e., only $\psi(x)$ or $\bar\psi(x)$ is needed. This is different on the
lattice. Both chiral rotation (7.31) and projection (7.38) require the
application of the lattice Dirac operator $D$. Thus the chiral rotation and
the decomposition into components involve neighboring sites and depend on the
gauge field. Consequently the chirality of a lattice fermion is determined
using information from the gauge field and from neighboring lattice sites.

\section{Consequences of the Ginsparg--Wilson equation}

In this section we discuss some of the most important consequences of the
Ginsparg--Wilson equation (7.29) without specifying a particular solution.
Actual solutions will be presented in Sect.~7.4, the so-called overlap action,
and in Sects.~9.2 and 9.3 fixed point fermions will be discussed. We remark
that chiral symmetry will also play a major role when we discuss Ward
identities in Chap.~11. This subject is not addressed here.

\subsection{Spectrum of the Dirac operator}

We begin our discussion of the consequences of the Ginsparg--Wilson equation
by analyzing the eigenvalue spectrum of the Dirac operator on a finite lattice.
The corresponding eigenvalue equation is denoted by
\begin{equation}
  D\,v_\lambda = \lambda\,v_\lambda .
\end{equation}
At the moment we do not invoke the Ginsparg--Wilson equation and only require
the Dirac operator $D$ to be $\gamma_5$-hermitian (compare Sect.~5.4.3),
\begin{equation}
  \gamma_5 D\gamma_5 = D^\dagger .
\end{equation}
This property is obeyed by most commonly used Dirac operators.\footnote{Unless
a chemical potential, a $\theta$-angle or a twisted mass term is introduced.}
For the example of the Wilson Dirac operator this can be shown in a few lines
of algebra, and the overlap operator, which we will discuss below, inherits
$\gamma_5$-hermiticity from the kernel Dirac operator used for its
construction.

The requirement of $\gamma_5$-hermiticity alone already has interesting
consequences for the eigensystem. For the characteristic polynomial $P(\lambda)$
of $D$ one finds (the asterisk denotes complex conjugation)
\begin{align}
  P(\lambda) &= \det[D - \lambda\,\mathbb{1}]
  = \det[\gamma_5^2(D-\lambda\,\mathbb{1})]
  = \det[\gamma_5(D-\lambda\,\mathbb{1})\gamma_5] \notag\\
  &= \det[D^\dagger - \lambda\,\mathbb{1}]
  = \det[D - \lambda^*\,\mathbb{1}]^*
  = P(\lambda^*)^*,
\end{align}
where we have used $\gamma_5^2 = 1$ and (7.44). The eigenvalues $\lambda$ are
the zeros of $P(\lambda)$ and (7.45) implies that if $\lambda$ is a zero, so is
$\lambda^*$. Thus for a $\gamma_5$-hermitian Dirac operator the eigenvalues
are either real or come in complex conjugate pairs.

Equation (7.44) also has an interesting consequence for the $\gamma_5$ matrix
element of the eigenvectors. Writing the inner product of two vectors $u$, $v$
as $u^\dagger v = (u, v)$ we find
\begin{align}
  \lambda\,(v_\lambda, \gamma_5 v_\lambda)
  &= (v_\lambda, \gamma_5 D v_\lambda)
  = (v_\lambda, D^\dagger\gamma_5 v_\lambda)
  = (D v_\lambda, \gamma_5 v_\lambda) \notag\\
  &= \lambda^*\,(v_\lambda, \gamma_5 v_\lambda),
\end{align}
and thus $(\mathrm{Im}\,\lambda)\,(v_\lambda, \gamma_5 v_\lambda) = 0$. This
implies
\begin{equation}
  (v_\lambda, \gamma_5 v_\lambda) = 0, \qquad \text{unless} \quad \lambda \in \mathbb{R} .
\end{equation}
Hence only eigenvectors $v_r$ with real eigenvalues $r$ can have nonvanishing
chirality, i.e., $(v_r, \gamma_5 v_r) \neq 0$.

In addition to being $\gamma_5$-hermitian, we now also require the Dirac
operator $D$ to obey the Ginsparg--Wilson equation (7.29). Multiplying this
equation with $\gamma_5$ from either left or right and using (7.44) we obtain
\begin{equation}
  D^\dagger + D = a\,D^\dagger D, \qquad
  D + D^\dagger = a\,D D^\dagger .
\end{equation}
A first consequence of these two equations is that a $\gamma_5$-hermitian
Ginsparg--Wilson Dirac operator $D$ is a normal operator, i.e., $D$ and
$D^\dagger$ commute, since the left-hand sides of the two Equations (7.48) are
equal. Normality implies that the eigenvectors form an orthogonal basis (see
Appendix~A.5). Furthermore, multiplying the first equation in (7.48) with a
normalized eigenvector $v_\lambda$ from the right and with $v_\lambda^\dagger$
from the left, one finds
\begin{equation}
  \lambda^* + \lambda = a\,\lambda^*\lambda .
\end{equation}
Writing the eigenvalue as $\lambda = x + iy$ this equation turns into
\begin{equation}
  \left(x-\frac{1}{a}\right)^2 + y^2 = \frac{1}{a^2},
\end{equation}
which shows that the eigenvalues of a $\gamma_5$-hermitian Ginsparg--Wilson
Dirac operator are restricted to a circle in the complex plane (compare the
left-hand side plot of Fig.~7.1). This so-called Ginsparg--Wilson circle has
its center at $1/a$ on the real axis and a radius of $1/a$. We remark that the
doubler modes, i.e., those modes where in the free case at least one of the
momentum components equals to $\pi/a$, end up near $2/a$ in the complex plane
and decouple as $a \to 0$.

A convenient parametrization of the eigenvalues $\lambda$ of a
Ginsparg--Wilson Dirac operator is given by
\begin{equation}
  \lambda = \frac{1}{a}\left(1 - e^{i\phi}\right), \qquad
  \phi \in (-\pi, \pi] .
\end{equation}
Using this parametrization one finds that the eigenvalues $1/\lambda$ of the
quark propagator $D^{-1}$ fall on a line parallel to the imaginary axis (see
Fig.~7.1),
\begin{equation}
  \frac{1}{\lambda} = \frac{a}{2} + i\,\frac{a\sin(\phi)}{2\,(1-\cos(\phi))} .
\end{equation}

\begin{figure}[htbp]
  \centering
  \includegraphics[width=0.9\textwidth]{fig71.png}
  \caption*{Fig. 7.1. \emph{Allowed regions for the eigenvalues $\lambda$ of a
  Ginsparg--Wilson Dirac operator in the complex plane (left-hand side plot)
  and for the eigenvalues $\lambda^{-1}$ of the corresponding propagator
  (right-hand side plot)}.}
  \label{fig:71}
\end{figure}

Since the Ginsparg--Wilson circle touches the origin of the complex plane,
$D$ may have exact zero eigenvalues. Let $v_0$ be a zero mode, then
\begin{equation}
  D v_0 = 0 \;\Rightarrow\; \gamma_5 D v_0 = 0 \;\Rightarrow\; D\gamma_5 v_0 = 0,
\end{equation}
using the Ginsparg--Wilson equation in the last step. Thus, on the subspace of
zero modes, the Dirac operator commutes with $\gamma_5$ and the zero modes can
be chosen as eigenstates of $\gamma_5$. Since $\gamma_5^2 = 1$, the eigenvalues
of $\gamma_5$ are $+1$ and $-1$ and we conclude
\begin{equation}
  \gamma_5 v_0 = \pm v_0,
\end{equation}
implying that the zero modes are chiral. A zero mode with positive chirality
is referred to as right handed, while left handed is used for negative
chirality.

From the left-hand side plot in Fig.~7.1 it is obvious that a Ginsparg--Wilson
Dirac operator can have real eigenvalues also at $2/a$. Following the same
steps as for the zero modes, one finds that also the eigenmodes with real
eigenvalue $2/a$ are chiral. These eigenmodes are the doubler partners of the
zero modes shifted to the other side of the Ginsparg--Wilson circle, where
they decouple in the limit $a \to 0$.

\subsection{Index theorem}

In the continuum the celebrated Atiyah--Singer index theorem [14] relates the
numbers of left- and right-handed zero modes of the massless Dirac operator to
a quantity $Q_{\mathrm{top}}$ which is a property of the gauge fields, the
so-called topological charge. To derive the lattice equivalent of the index
theorem [15] we consider
\begin{align}
  Q_{\mathrm{top}} &\equiv \frac{1}{2}\tr[\gamma_5 D]
  = -\frac{1}{2}\tr[\gamma_5(2-aD)]
  = -\frac{1}{2}\sum_\lambda (v_\lambda, \gamma_5(2-aD)v_\lambda) \notag\\
  &= -\frac{1}{2}\sum_\lambda (2-a\lambda)\,(v_\lambda, \gamma_5 v_\lambda)
  = n_- - n_+,
\end{align}
where $n_-$, $n_+$ denote the numbers of left- and right-handed zero modes.
In the second step we have used $\tr[\gamma_5] = 0$, and in the third step we
expressed the trace as a sum over the eigenvectors $v_\lambda$ of $D$. As
discussed above, $D$ is normal and thus its normalized eigenvectors form an
orthonormal basis suitable for computing the trace. In the last step we have
used (7.47), implying that only eigenvectors with real eigenvalues have
nonvanishing chirality and that the factor $(2-a\lambda)$ cancels the
contributions from the doubler modes with $\lambda = 2/a$. The $\gamma_5$
matrix elements of the surviving zero modes are $\pm 1$ due to (7.54).

If one leaves out the second step in (7.55), where the extra $\gamma_5$ was
sneaked in, one obtains $Q_{\mathrm{top}} = n'_+ - n'_-$, where $n'_\pm$ are
the numbers of left- and right-handed eigenvectors with eigenvalue $2/a$. This
gives the same result, since the doublers enter with opposite chirality.

The topological charge $Q_{\mathrm{top}}$ can be written as the space--time
sum of the so-called topological charge density $q(n)$. Doing so we obtain
\begin{equation}
  Q_{\mathrm{top}} = a^4\sum_{n\in\Lambda} q(n), \qquad
  q(n) = \frac{1}{2a^3}\,\tr_{CD}[\gamma_5 D(n|n)], \qquad
  Q_{\mathrm{top}} = n_- - n_+,
\end{equation}
where $\tr_{CD}$ denotes the trace over color and Dirac indices.

The lattice index theorem (7.56) is a truly remarkable result: The Dirac
operator depends on the gauge field variables $U_\mu(n)$. Thus the topological
charge $Q_{\mathrm{top}}$ is a functional of the gauge field, and it is a
highly nontrivial statement that this functional is an integer number. As the
name ``topological charge'' already indicates the underlying reason is of
topological nature.

At this point we remark that in actual lattice calculations one never finds
zero modes of both chiralities for a single configuration, i.e., $n_+$ and
$n_-$ are not both nonzero simultaneously. It may be argued that the
underlying reason is that only a very subtle arrangement of the gauge field
would give rise to $n_+ \neq 0$ and $n_- \neq 0$ at the same time. Such an
arrangement has zero measure in the path integral, and the fact that at least
one of the numbers $n_+$, $n_-$ vanishes implies $Q_{\mathrm{top}} = n_-$ or
$Q_{\mathrm{top}} = -n_+$. This fact is sometimes referred to as absence of
fine tuning.

In the continuum the index theorem takes the form
$Q^{\mathrm{cont}}_{\mathrm{top}} = n_- - n_+$, with the topological charge
given by
\begin{equation}
  Q^{\mathrm{cont}}_{\mathrm{top}} = \int d^4x\;q(x)^{\mathrm{cont}}, \qquad
  q(x)^{\mathrm{cont}} = \frac{1}{32\pi^2}\,\varepsilon_{\mu\nu\rho\sigma}
  \tr_C[F_{\mu\nu}(x)F_{\rho\sigma}(x)],
\end{equation}
where $\varepsilon_{\mu\nu\rho\sigma}$ is the completely anti-symmetric rank 4
Levi--Civita tensor and $F_{\mu\nu}$ is the field strength tensor (2.15). The
functional $Q^{\mathrm{cont}}_{\mathrm{top}}$ is an expression for the
so-called Pontryagin index.

An instance of a continuum gauge field configuration with nonvanishing
$Q^{\mathrm{cont}}_{\mathrm{top}}$ is so-called instantons [16]. There it was
also shown that instantons are local minima of the gauge action. Thus they
give a contribution in a saddle point evaluation of the path integral. This
insight led to a wealth of papers analyzing the role of topological objects in
QCD (for reviews see [17, 18] and references therein). For an introductory
text on topological field configurations and an elementary proof that
$Q^{\mathrm{cont}}_{\mathrm{top}}$ of (7.57) is indeed integer valued we
recommend [19].

On the lattice the topological charge (7.55) is an integer due to its
definition through the number of zero modes of a chiral Dirac operator. Thus
the remaining question is whether the lattice definition (7.56) of the
topological charge density $q(n)$ approaches its continuum counterpart
$q(x)^{\mathrm{cont}}$ of (7.57) in the limit $a \to 0$. It can be shown that
under certain (smoothness) conditions for the gauge fields one has
$q(n) = q(x)^{\mathrm{cont}} + \order{a^2}$ as $a \to 0$. Several derivations
of this fact have been given for various settings and conditions on the gauge
fields, the first results being [15, 20, 21].

The QCD path integral is symmetric with respect to configurations of positive
and negative topological charges and thus the expectation value
$\ev{Q_{\mathrm{top}}}$ vanishes. However, one can consider the topological
susceptibility defined as
\begin{equation}
  \chi_{\mathrm{top}} = \frac{1}{V}\ev{Q_{\mathrm{top}}^2}
  = \frac{1}{a^8|\Lambda|}\sum_{m,n}\ev{q(m)q(n)}
  = a^4\sum_n \ev{q(0)q(n)},
\end{equation}
where in the last step we have used translational invariance of the correlator
to get rid of the summation over $m$. We remark that $\chi_{\mathrm{top}}$ is
volume dependent. The topological susceptibility considered in the infinite
volume limit provides information about the distribution of topological charge
as a function of $\Nf$ and $m$ [22]. Using the fermionic definition (7.55) of
the topological charge it has been determined in several quenched calculations
[23--25] (and references therein) with a result of
$\chi_{\mathrm{top}} \approx (190\,\mathrm{MeV})^4$. Since in the quenched case
the fermion determinant is not taken into account, this result is independent
of $m$. For the dynamical case the simulations are not so far advanced and the
dependence of $\chi_{\mathrm{top}}$ on $\Nf$ and $m$ is not finally settled.
Phenomenologically, the topological susceptibility plays an interesting role
in the so-called Witten--Veneziano formula which we briefly address at the end
of the next section.

\subsection{The axial anomaly}

Having implemented chiral symmetry on the lattice, we are now ready to give a
(lattice) derivation [12] of the axial anomaly which we have already addressed
in Sect.~7.1. There it was stated that the integration measure for the fermions
is not invariant under the flavor diagonal rotation (7.16).

To derive the anomaly on the lattice we consider lattice QCD with $\Nf$
massless flavors with a Dirac operator obeying the Ginsparg--Wilson equation.
The spinors $\psi$, $\bar\psi$, describing our $\Nf$ flavors, are transformed
with an infinitesimal chiral rotation of the form
\begin{equation}
  \psi' = \left(1 + i\varepsilon M\gamma_5\left(1-\frac{aD}{2}\right)\right)\psi,
  \qquad
  \bar\psi' = \bar\psi\left(1 + i\varepsilon M\left(1-\frac{aD}{2}\right)\gamma_5\right).
\end{equation}
This transformation is the infinitesimal version of (7.31) obtained by setting
$\alpha = \varepsilon$ in (7.31) and keeping only the leading terms in the
expansion of the exponential in $\varepsilon$. In addition we consider several
flavors by allowing for a mixing matrix $M$ which can either be the identity,
$M = \mathbb{1}_{\Nf}$, or one of the generators of $SU(\Nf)$, $M = T_i$.
According to (5.26) for the transformation of the integration measure of
Grassmann variables we obtain
\begin{align}
  \mathcal{D}[\psi,\bar\psi] &= \mathcal{D}[\psi',\bar\psi']\,
  \det\!\left[1 + i\varepsilon M\gamma_5\left(1-\frac{aD}{2}\right)\right]
  \det\!\left[1 + i\varepsilon M\left(1-\frac{aD}{2}\right)\gamma_5\right] \notag\\
  &= \mathcal{D}[\psi',\bar\psi']\,
  \det\!\left[1 + i\varepsilon M\gamma_5\left(1-\frac{aD}{2}\right)\right]^2,
\end{align}
where in the second step we used the invariance of the determinant under cyclic
permutation of $\gamma_5$. Expansion in $\varepsilon$ gives
\begin{align}
  \det\Bigl[1 + i\varepsilon M\gamma_5\Bigl(1-\tfrac{aD}{2}\Bigr)\Bigr]^2
  &= \exp\Bigl[2\tr\ln\Bigl(1 + i\varepsilon M\gamma_5\Bigl(1-\tfrac{aD}{2}\Bigr)\Bigr)\Bigr] \notag\\
  &= \exp\Bigl[-2\sum_{j=1}^\infty \frac{(-i\varepsilon)^j}{j}\,
     \tr M\Bigl(\gamma_5\Bigl(1-\tfrac{aD}{2}\Bigr)\Bigr)^j\Bigr] \notag\\
  &= 1 + 2i\varepsilon\,\tr\Bigl[M\gamma_5\Bigl(1-\tfrac{aD}{2}\Bigr)\Bigr] + \order{\varepsilon^2} \notag\\
  &= 1 + 2i\varepsilon\,\tr_F[M]\sum_{n\in\Lambda}\tfrac{a}{2}\Bigl(3\,\tr_D[\gamma_5] \notag\\
  &\qquad - \tr_{CD}[\gamma_5 D(n|n)]\Bigr) + \order{\varepsilon^2}.
\end{align}
In the first line of this equation we have used a well-known formula for the
determinant (see Appendix~A.5) and in the second line the logarithm was
expanded in a power series. In the last line we have split the trace into its
components, with $\tr_F$, $\tr_C$, and $\tr_D$ denoting the traces over flavor,
color, and Dirac indices, respectively. The sum runs over all sites of the
lattice $\Lambda$. Since $\gamma_5$ is traceless, the first of the
$\order{\varepsilon}$ terms vanishes.

If the flavor matrix $M$ is chosen to be one of the $SU(\Nf)$ generators, the
anomaly is canceled due to the vanishing trace over the flavor indices. Only
for the flavor singlet choice $M = \mathbb{1}_{\Nf}$ this trace is nontrivial,
$\tr_F[M] = \Nf$. Combining (7.60) and (7.61) we obtain for this case
\begin{equation}
  \mathcal{D}[\psi,\bar\psi]
  = \mathcal{D}[\psi',\bar\psi']\left(1 - 2i\varepsilon\,\Nf\,Q_{\mathrm{top}}
     + \order{\varepsilon^2}\right),
\end{equation}
where we have used definition (7.56) of the topological charge
$Q_{\mathrm{top}}$.

We conclude that indeed a chiral, flavor singlet rotation is not a symmetry of
QCD, although the action itself is invariant (for massless quarks). The
non-invariance comes from the measure and the symmetry breaking term is the
topological charge. This result underlines again the strong connection between
chiral symmetry and topological gauge field configurations which we have
already addressed when we discussed the index theorem.

The axial anomaly has important physical consequences. Most importantly it
implies that the chiral flavor singlet symmetry cannot be broken spontaneously
and no flavor singlet Goldstone particle exists. This explains why the
$\eta$-meson is much heavier than the pions, when one considers approximate
chiral symmetry for $\Nf = 2$. For approximate chiral symmetry with $\Nf = 3$,
this role is taken over by the $\eta'$-meson. The aforementioned
Witten--Veneziano formula [26--28] relates the mass of the $\eta'$ to the
topological susceptibility $\chi_{\mathrm{top}}$, thus revealing a deep
connection of QCD vacuum properties and the mass spectrum. The mathematical
subtleties of the Witten--Veneziano formula have been addressed in several
recent articles [29--33], partly based on the rigorous framework provided by
Ginsparg--Wilson fermions.

\subsection{The chiral condensate}

In Sect.~7.1 we have discussed that in the continuum the vacuum expectation
value $\ev{u(x)\bar{u}(x)}$ is not invariant under any of the chiral rotations
(7.15) and (7.16). Thus it can serve as an order parameter for spontaneous
chiral symmetry breaking, with a nonvanishing condensate signaling the
spontaneous breaking of chiral symmetry. We emphasize again that this
spontaneous breaking occurs only for $\Nf \geq 2$, since for $\Nf = 1$ the
single axial symmetry is not broken spontaneously, but explicitly by the
anomaly.

When discussing chiral symmetry we have identified the correct mass term on
the lattice in (7.40), i.e., the scalar bilinear which maximally mixes left-
and right-handed components. Accordingly, we consider the scalar expectation
value for lattice QCD with Ginsparg--Wilson fermions
\begin{equation}
  \Sigma^{\mathrm{lat}}(a, m, |\Lambda|) \equiv
  -\left\langle \bar{u}(n)\left(1-\frac{aD}{2}\right)u(n)\right\rangle .
\end{equation}
The expectation value on the right-hand side is evaluated on a lattice
$\Lambda$ with $|\Lambda|$ sites and a lattice spacing $a$. We consider $\Nf$
flavors of fermions, all with mass $m$, and the corresponding Dirac operator
$D_m$ is given by (7.41). The term $aD/2$ in (7.63) cancels the real part of
the eigenvalues (7.52) of the massless propagator, thus shifting its spectrum
onto the imaginary axis as in the continuum (see the right-hand side of
Fig.~7.1).

Since our formulation is translationally invariant, (7.63) is independent of
$n$ and we can average over all lattice sites. Doing so and performing the
Grassmann integration according to (5.54) we find
\begin{equation}
  \Sigma^{\mathrm{lat}}(a, m, |\Lambda|) =
  \frac{1}{a^4|\Lambda|}\left\langle
    \tr\left[\left(1-\frac{aD}{2}\right)D_m^{-1}\right]
  \right\rangle_G,
\end{equation}
where the trace runs over space--time, color, and Dirac indices, and
\begin{equation}
  \ev{X}_G = \frac{1}{Z}\int \mathcal{D}[U]\,e^{-S_G[U]}\,\det[D_m]^{\Nf}\,X
\end{equation}
is the remaining gauge field integral. With the abbreviation (7.42) the
propagator, i.e., the inverse of $D_m$ defined in (7.41), reads
$D_m^{-1} = (\omega D + m\,\mathbb{1})^{-1}$, $\omega = 1 - am/2$, and after
some algebra (7.64) turns into
\begin{equation}
  \Sigma^{\mathrm{lat}}(a, m, |\Lambda|) =
  \frac{1}{\omega\, a^4|\Lambda|}\left\langle
    \tr\left[(\omega D + m\,\mathbb{1})^{-1}\left(1-\frac{aD}{2}\right)\right]
  \right\rangle_G .
\end{equation}
This expression is now ready for a numerical or analytical investigation of
the chiral condensate.

At this point a few remarks on the necessary limits are in order. Spontaneous
breaking of a symmetry requires an infinite system, therefore the limit
$|\Lambda| \to \infty$ has to be taken. Subsequently the explicit breaking
through the mass term has to be removed, i.e., the limit $m \to 0$ has to be
performed. Like in other cases of spontaneous symmetry breaking these limits
cannot be interchanged. After these limits, (7.66) becomes a bare condensate at
a fixed ultraviolet cutoff $1/a$. The definition of the physical condensate is
regularization scheme dependent and its value must be converted to other
schemes (like the continuum $\overline{MS}$ scheme) by renormalization with
the scalar renormalization factor $Z_S$ which we discuss in Chap.~11.
Subsequently one can analyze the scaling with $a$ and perform the continuum
limit.

It is obvious that in a numerical approach the discussed sequence of limits is
hard to establish. However, in [22] the chiral condensate has been studied in
a 4D box of volume $V$ at finite mass $m$ using low-energy effective theory.
An expression for the chiral condensate $\Sigma(m, V)$ as a function of $m$ and
$V$ was obtained which contains the physical chiral condensate
\begin{equation}
  \Sigma \equiv \lim_{m \to 0}\;\lim_{V \to \infty} \Sigma(m, V)
\end{equation}
as a parameter. Here both, $\Sigma$ and $\Sigma(m, V)$, are understood at
finite cutoff. Since this is not necessarily a lattice cutoff we do not
indicate it explicitly. Once the functional form $\Sigma(m, V)$ is known, one
can use lattice simulations to compute the condensate
$\Sigma^{\mathrm{lat}}(a, m, |\Lambda|)$ at finite volume $V = a^4|\Lambda|$
and finite mass $m$ and fit the data to the function $\Sigma(m, V)$ to obtain
the true bare condensate $\Sigma$ as a fit parameter. This procedure has first
been carried out for the quenched case, which we discuss here. As always, the
motivation for studying the quenched case is the hope that the quenched result
for $\Sigma$ will be close to its value at $\Nf = 3$.

Within partially quenched chiral perturbation theory one can derive the
functional form for the chiral condensate [34, 35] for different topological
sectors as a function of the volume $V$ and mass parameter $m$. We speak of the
topological sector $\nu$ when in the path integral only gauge configurations
with topological charge $Q_{\mathrm{top}} = \nu$ are considered. On the lattice
this is a concept which can, e.g., be implemented by sorting all configurations
of a Monte Carlo ensemble with respect to $Q_{\mathrm{top}}$ as determined from
the index theorem (7.56). The result from chiral perturbation theory mirrors a
finding from random matrix theory, where one obtains universal properties of
the spectra of certain symmetry classes of random matrices. The quenched result
for the condensate in the topological sector $\nu$ is given by
\begin{equation}
  \Sigma(m, V)_\nu = \Sigma\left[
    z\Bigl(I_{|\nu|}(z)K_{|\nu|}(z) + I_{|\nu|+1}(z)K_{|\nu|-1}(z)\Bigr)
    + \frac{|\nu|}{mV}\right],
\end{equation}
where $z = m V \Sigma$ is a dimensionless scaling variable and $\Sigma$ is the
bare condensate defined in (7.67). The lattice condensate
$\Sigma^{\mathrm{lat}}(a, m, |\Lambda|)$ is then computed for each sector $\nu$
separately and fitted to the functional form (7.68).

Let us discuss the strategy for a numerical determination of $\Sigma$ in more
detail [36]. Expanding the Bessel functions in (7.68) one finds for $m \to 0$:
\begin{figure}[htbp]
  \centering
  \includegraphics[width=0.9\textwidth]{fig72.png}
  \caption*{Fig. 7.2. \emph{Lattice results for the bare condensate for
  different volumes and different topological sectors [37]}. (Reprinted from
  Hasenfratz et al., copyright (2002) with permission from Elsevier)}
  \label{fig:72}
\end{figure}
\begin{equation}
  \frac{\Sigma(m, V)_\nu}{m}
  \xrightarrow{m \to 0}
  \begin{cases}
    \displaystyle \Sigma^2 V\left(\frac{1}{2} - \ln(z/2) - \gamma\right),
      & \nu = 0,\\[3mm]
    \displaystyle \frac{\Sigma^2 V}{2|\nu|} + \frac{|\nu|}{2\,mV},
      & \nu \neq 0 .
  \end{cases}
\end{equation}
The topologically trivial case ($\nu = 0$) contains a logarithmic divergence at
small $m$ ($=$ small $z$) and thus was found to be not very suitable for a
numerical study at small $m$. For $\nu \neq 0$ an extra term from the zero
modes appears which diverges even more strongly.\footnote{For the case of
dynamical fermions, zero modes are suppressed by the fermion determinant which
vanishes at $m = 0$ and thus no such divergent term occurs.} However, this
contribution can be cleanly removed by hand in order to improve the signal.
This is done by exploiting the chirality of the zero modes (7.54) and taking
the trace in (7.66) only after the propagator has been multiplied with
$(1 \pm \gamma_5)/2$. For each individual gauge configuration the sign is
chosen opposite to the chirality of the zero modes. After this modification
the trace can be computed using the stochastic estimator techniques discussed
in Chap.~8.

Since we are computing the bare condensate at finite $V$ and $m$, cutoff
effects can alter behavior (7.69). In leading order in $m$ this gives rise to
an extra cutoff-dependent constant $c$. Thus for nontrivial topological sectors
and subtracted zero mode contributions one finds the final form
\begin{equation}
  \frac{\Sigma(m, V)'_\nu}{m}
  \xrightarrow{m \to 0}
  \frac{\Sigma^2 V}{2|\nu|} + c
  \qquad (\nu \neq 0) .
\end{equation}
The ratio $\Sigma^{\mathrm{lat}}(a, m, |\Lambda|)/m$ thus is expected to
approach a constant at $m \to 0$, and comparison of different volumes and
topological sectors allows one to extract the bare condensate $\Sigma$. The
data in Fig.~7.2 were computed in [37] along these lines and nicely illustrate
the behavior according to (7.70).

The last step, which we will address in detail only in Chap.~11, is the
renormalization with $Z_S$ to obtain the physical quenched condensate
$\Sigma_{\mathrm{phys}} = Z_S \Sigma$. Several independent calculations
[36--39] determined the quenched chiral condensate and after renormalization
obtained $\Sigma_{\mathrm{phys}} \approx (270\,\mathrm{MeV})^3$ (in the
$\overline{MS}$ scheme at renormalization scale 2\,GeV).

\subsection{The Banks--Casher relation}

So far we have discussed the numerical analysis of the chiral condensate for a
finite system at finite quark mass $m$. However, one can at least formally
perform the thermodynamic and chiral limits and parameterize the unknown
nonperturbative information in terms of the eigenvalue density of the Dirac
operator near the origin. This gives rise to an important relation, the
so-called Banks--Casher relation [40], which we now derive.

The starting point of our derivation is (7.66). In a first step we express the
trace as a sum over all eigenvalues $\lambda_i$ of $D$:
\begin{align}
  \tr\left[(\omega D + m\,\mathbb{1})^{-1}\left(1-\frac{aD}{2}\right)\right]
  &= \sum_{\lambda_i}\left(\frac{1}{\omega\lambda_i + m} - \frac{a}{2}\right) \notag\\
  &= (n_+ + n_-)\left(\frac{1}{m} - \frac{a}{2}\right) \notag\\
  &\quad + (n'_+ + n'_-)\left(\frac{1}{2\omega/a + m} - \frac{a}{2}\right) \notag\\
  &\qquad + \sum_{\lambda_i \neq 0,\,2/a}\left(\frac{1}{\omega\lambda_i + m} - \frac{a}{2}\right),
\end{align}
where in the second line we have split the sum over all eigenvalues in the
contributions of the $n_+ + n_-$ eigenvalues $0$, the $n'_+ + n'_-$ eigenvalues
$2/a$, and a sum over the complex eigenvalues. Recalling definition (7.42) for
$\omega$ one finds that the factor for the contribution of the eigenvalues at
$0$ reduces to $\omega/m$, whereas the contribution of the eigenvalues at
$2/a$ vanishes. Expressing the eigenvalues $\lambda_i$ through the phases
$\phi_i$ introduced in (7.51) and combining the complex conjugate pairs of
eigenvalues the last expression turns into
\begin{equation}
  \frac{\omega}{m}(n_+ + n_-) + \sum_{\phi_i \neq 0,\pi}
  \frac{\omega\,(1+\cos\phi_i)\,m}{(2-2\cos\phi_i)(1-a^2m^2/4)/a^2 + m^2} .
\end{equation}
Thanks to the above discussed ``absence of fine tuning'' for a given gauge
configuration only one of the numbers $n_+$, $n_-$ can be nonvanishing and we
may replace $(n_+ + n_-) \to |Q_{\mathrm{top}}|$, the number of zero modes.
Inserting (7.72) for the trace in (7.66) one arrives at
\begin{align}
  \Sigma^{\mathrm{lat}}(a) &\equiv \lim_{m\to 0}\;\lim_{|\Lambda|\to\infty}
  \Sigma^{\mathrm{lat}}(a,m,|\Lambda|) \notag\\
  &= \lim_{|\Lambda|\to\infty}\frac{1}{a^4|\Lambda|}\,\frac{\ev{|Q_{\mathrm{top}}|}_G}{m}
   \notag\\
  &\qquad + \frac{1}{2a^4}\int_{-\pi}^{\pi}d\phi\,\rho_A(\phi)\,
     \frac{(1+\cos\phi)\,m}{(2-2\cos\phi)(1-a^2m^2/4)/a^2 + m^2},
\end{align}
where we have defined the angular density $\rho_A(\phi)$ for the angles
$\phi$ of the eigenvalues on the Ginsparg--Wilson circle as
\begin{equation}
  \rho_A(\phi) = \left\langle \lim_{|\Lambda|\to\infty}\frac{1}{|\Lambda|}
  \sum_{\phi_i \neq 0,\pi}\delta(\phi-\phi_i)\right\rangle_G .
\end{equation}
In the definition of the density $\rho_A(\phi)$ we have already performed the
limit $|\Lambda| \to \infty$ (at fixed $a$) necessary for the spontaneous
breaking of a continuous symmetry. In this limit the eigenvalues become dense
on the Ginsparg--Wilson circle and $\rho_A(\phi)$ as defined in (7.74) indeed
becomes a density.

The first contribution (7.73) vanishes for $|\Lambda| \to \infty$, since
$\ev{|Q_{\mathrm{top}}|}_G$ does not grow faster than $\sqrt{V}$. This follows
from the finiteness of $\chi_{\mathrm{top}} = \ev{Q^2_{\mathrm{top}}}/V$.
For the second term in (7.73) one can utilize that
\begin{equation}
  \delta_m(X) = \frac{1}{\pi}\,\frac{m}{X^2 + m^2}\left(1 + \order{m^2}\right)
\end{equation}
is a $\delta$-sequence, i.e., $\delta_m(X) \to \delta(X)$ for $m \to 0$ (see
e.g. [41]). Thus we obtain for the second part of (7.73), which remains finite
for $m \to 0$:
\begin{align}
  &\frac{1}{2a^4}\int_{-\pi}^{\pi}d\phi\,\rho_A(\phi)\,
    \frac{(1+\cos\phi)\,m}{(2-2\cos\phi)(1-a^2m^2/4)/a^2 + m^2} \notag\\
  &= \frac{\pi}{2a^3}\int_{-\pi}^{\pi}d\phi\,(1+\cos\phi)\,
     \delta\!\bigl(a^{-1}(2-2\cos\phi)\bigr)\,\rho_A(\phi) \notag\\
  &= \frac{\pi}{2a^3}\int_{-\pi}^{\pi}d\phi\,(1+\cos\phi)\,
     \delta(\phi)\,\rho_A(\phi) = \frac{\pi}{a^3}\,\rho_A(0),
\end{align}
where we have used standard manipulations for the Dirac-delta with a function
as argument. Equation (7.76) shows that after taking the infinite volume limit
the chiral condensate is proportional to the angular density $\rho_A(\phi)$ at
$\phi = 0$. We emphasize that in the definition of the density the exact zero
eigenvalues are left out.

Usually one does not use the angular density $\rho_A$, but the density
$\rho_\lambda$ of eigenvalues on the imaginary axis is
$\rho_\lambda = \Delta n/\Delta y$, where $\Delta n$ is the number of
eigenvalues per interval $\Delta y$ on the imaginary axis. For small angles
$\phi$ one has $\Delta y = \Delta\phi/a$ (compare Fig.~7.1), implying that
$\rho_\lambda = a\,\rho_A$. Thus we conclude
\begin{equation}
  \frac{\pi}{a^3}\,\rho_A(0) = \frac{\pi}{a^4}\,\rho_\lambda(0)
  = \pi\,\rho(0),
\end{equation}
where $\rho$ is the density of eigenvalues on the imaginary axis per unit
volume.

While the numerical calculation presented in the last section shows the
existence of a nonvanishing chiral condensate from first principles, it cannot
provide insight about an underlying physical mechanism for chiral symmetry
breaking. The Banks--Casher relation (7.77), however, has at least opened the
door to understanding such a mechanism.

A very influential idea is that the chiral condensate is formed through a
``fluid'' of weakly interacting topological field configurations such as
instantons (see [17, 18] for reviews). Often, as in the case of instantons,
the topological configurations are localized and such structures have been
nicknamed topological lumps. A gauge field configuration which carries
topological charge gives rise to zero eigenvalues according to the index
theorem (7.56). For a mixture of topological lumps with different charges one
cannot assume that the zero eigenvalues survive unperturbed. Instead one
expects that the eigenvalues move in the imaginary direction. A fluid of
weakly interacting topological lumps gives rise to an accumulation of
eigenvalues on the imaginary axis thus building up the density of eigenvalues
$\rho(0)$ near the origin. This in turn gives rise to the nonvanishing chiral
condensate via the Banks--Casher relation (7.77). This mechanism and the
possible structure of the topological lumps have been studied extensively on
the lattice and the reviews [42--46] may serve as a guide to further reading.

The density of eigenvalues can be studied in some detail within random matrix
theory, which we addressed briefly in the preceding section. Fine details like
distribution densities of the smallest eigenvalue, the next-to-smallest
eigenvalue, and so on appear to be universal properties of the spectra of
random matrices with the symmetry properties of the Dirac operator. There is a
wealth of tools to analyze spectra in this way and thereby disentangle
``universal'' features from properties specific for the dynamics of the
theory, such as the condensate value. More about these issues can be found in
[47, 48].

\section{The overlap operator}

Up to now our presentation of chiral symmetry on the lattice did not refer to
a particular solution of the Ginsparg--Wilson equation. In this section we
close this gap and introduce the overlap operator. We analyze the locality
properties of the overlap operator and discuss its numerical evaluation. A
second solution of the Ginsparg--Wilson equation, the fixed point Dirac
operator, will be presented in Chap.~9. Furthermore, in Sect.~10.2 we discuss
domain wall fermions which are closely related to the overlap formulation.

\subsection{Definition of the overlap operator}

Originating from the initial papers on the overlap formulation [49--52],
Neuberger presented the modern form of the overlap Dirac operator in [53] and
showed in [11] that it is a solution of the Ginsparg--Wilson equation (7.29).
Explicitly the overlap Dirac operator is given by
\begin{equation}
  D_{\mathrm{ov}} = \frac{1}{a}\left(1 + \gamma_5\,\mathrm{sign}[H]\right),
  \qquad
  H = \gamma_5 A,
\end{equation}
where $A$ denotes some suitable $\gamma_5$-hermitian ``kernel'' Dirac operator,
thus $H$ is hermitian with real eigenvalues. The operator sign function may be
defined through the spectral representation of the operator (A.51).
Alternatively one may write the overlap operator in the form
\begin{equation}
  D_{\mathrm{ov}} = \frac{1}{a}\left(1 + \gamma_5 H\left(H^2\right)^{-1/2}\right).
\end{equation}
The simplest choice is to use the Wilson Dirac operator for the kernel,
\begin{equation}
  A = a\,D_W - \mathbb{1}(1+s), \qquad
  \text{such that} \qquad
  D_{\mathrm{ov}} = \frac{1}{a}\left(1 + A\left(\gamma_5 A\gamma_5 A\right)^{-1/2}\right),
\end{equation}
where $s$ is a real parameter with $|s| < 1$ which can be used to optimize
locality (see below). $D_W$ is the massless Wilson Dirac operator (compare
(5.51))
\begin{equation}
  D_W(n|m)^{ab}_{\alpha\beta} = \frac{1}{a}\,\delta_{\alpha\beta}\,\delta_{ab}\,
  \delta_{n,m}
  - \frac{1}{2a}\sum_{\mu=\pm1}^{\pm4}(1-\gamma_\mu)_{\alpha\beta}\,
    U_\mu(n)_{ab}\,\delta_{n+\hat\mu,m} .
\end{equation}
Since the Wilson Dirac operator is $\gamma_5$-hermitian (compare (5.76)) so is
$A$, i.e., $\gamma_5 A\gamma_5 = A^\dagger$. Thus we have
$\gamma_5 A\gamma_5 A = A^\dagger A$ and the square root in (7.80) is well
defined through the spectral theorem.

In order to see that the overlap operator obeys the Ginsparg--Wilson equation
we insert $D_{\mathrm{ov}}$ into (7.29). With the representation (7.78) we
find
\begin{align}
  a\,D_{\mathrm{ov}}D_{\mathrm{ov}}^\dagger
  &= \frac{1}{a}\left(1+\gamma_5\,\mathrm{sign}[H]\right)
     \left(1+\mathrm{sign}[H]\,\gamma_5\right) \notag\\
  &= \frac{1}{a}\left(1+\gamma_5\,\mathrm{sign}[H] + \mathrm{sign}[H]\,\gamma_5 + 1\right)
  = D_{\mathrm{ov}} + D_{\mathrm{ov}}^\dagger .
\end{align}
We have taken into account hermiticity, $\mathrm{sign}\,H =
\mathrm{sign}\,H^\dagger$, and the squared sign function as the unit operator:
$(\mathrm{sign}\,H)^2 = 1$. Similar to the steps (2.40), (2.41), and (2.42)
performed for the naive Dirac operator, one can expand the overlap operator
(7.80) for small $a$ and show that it approaches the Dirac operator of the
continuum (up to an irrelevant multiplicative constant:
$\hat{D}^{\mathrm{ov}} \approx i\gamma_\mu p_\mu/(1+s) + \order{p^2}$).

The Wilson operator $D_W$ used in the construction of the overlap operator can
be replaced by some other lattice Dirac operator which is free of doublers
[54]. This can improve the locality properties of the overlap operator (see
next section) and it may speed up the convergence in the numerical evaluation
of the overlap operator (Sect.~7.4.3) [38]. In this sense the overlap
construction may also be viewed as a projection of a non-chiral lattice Dirac
operator onto a solution of the Ginsparg--Wilson equation.

\subsection{Locality properties of chiral Dirac operators}

While the Wilson Dirac operator (7.81) and thus $A$ involve only nearest
neighbor terms and therefore are sparse matrices, the inverse square root
$(\gamma_5 A\gamma_5 A)^{-1/2}$ leads to a matrix that has nonvanishing entries
for all pairs $n$, $m$ of lattice sites. Thus $D_{\mathrm{ov}}(n|m) \neq 0$ for
all $n$, $m$ and the overlap operator is not an ultralocal operator. General
arguments that this is true for all solutions of the Ginsparg--Wilson equation
were given in [55].

However, locality is an essential concept for quantum field theories and its
violation is punished by noncausal interactions. Thus one has to understand in
which sense locality has to be manifest for a lattice field theory. A natural
requirement is that the Dirac operator $D$ falls off exponentially independent
of $\beta$, i.e., $D$ obeys the bound
\begin{equation}
  \left|D(n|m)^{ab}_{\alpha\beta}\right| \leq C\,\exp\left(-\gamma\,\|n-m\|\right),
\end{equation}
where the constants $C$ and $\gamma$ are independent of the gauge field. If
such a bound holds, then the interaction range $1/\gamma$ is a fixed distance
in lattice units. However, when expressed in physical units, this distance
decreases as we let $a \to 0$ in order to approach the continuum limit. Thus,
the interaction range in physical units, $a/\gamma$, shrinks to zero and in
the continuum limit we recover a local field theory. For the overlap operator
locality in the sense of (7.83) was established in [56], both with analytical
arguments and a numerical investigation. There it was also demonstrated that
the parameter $s$ in (7.80) can be tuned to optimize locality.

\subsection{Numerical evaluation of the overlap operator}

Numerically the main problem is to compute the sign function of the operator
$H$. Formally it is well defined through the spectral theorem:
\begin{equation}
  \mathrm{sign}[H] = \mathrm{sign}\left[\sum_i \lambda_i \ket{i}\bra{i}\right]
  = \sum_i \mathrm{sign}(\lambda_i)\,\ket{i}\bra{i} .
\end{equation}
Only for exactly vanishing eigenvalues of $H$, a numerically highly improbable
case, one introduces a ``tie-breaker,'' e.g., a randomly chosen $\pm 1$
replacing sign $\lambda$. However, for the huge Dirac matrices exact complete
diagonalization is prohibitively expensive and therefore methods based on
(7.84) are not applied except for some few test cases.

In most computations one utilizes instead
\begin{equation}
  \mathrm{sign}[H] = H\,|H|^{-1} = H\left(H^2\right)^{-1/2}
\end{equation}
and approximates $(H^2)^{-1/2}$ by either a polynomial in $H^2$ or a ratio of
polynomials. The convergence of such an approximations depends on the actual
matrix $H$, in particular on its eigenvalues, and therefore also on the current
gauge field configuration. Smaller eigenvalues of $H$ will lead to worse
convergence in general. As discussed in Sect.~6.2.5 the fluctuations of the
real eigenvalues of the Wilson Dirac operator are nonnegligible. A real
eigenvalue of $D_W$ in the vicinity of $(1+s)/a$ will give rise to a small
eigenvalue of $H$ leading to numerical problems. The situation can be improved
by applying smearing techniques (compare Sect.~6.2.6) or by using for the
overlap projection a different kernel operator which already approximates a
solution of the Ginsparg--Wilson equation.

For the numerical evaluation of the sign function we discuss here the two most
popular approximation methods and the so-called small eigenmode reduction
technique. Comparative studies are, e.g., found in [57--59].

\textbf{Polynomial approximation:} The Chebyshev polynomials $T_n(x)$ are
orthogonal polynomials with regard to the scalar product
\begin{equation}
  (f, g) \equiv \int_{-1}^{1} dx\,\frac{f(x)g(x)}{\sqrt{1-x^2}},
\end{equation}
normalized such that $(T_n, T_m) = \delta_{nm}$. A given function $r(x)$ has
the series expansion
\begin{equation}
  r(x) = \sum_{n=0}^{\infty} c_n\,T_n(x) \qquad \text{with} \qquad
  c_n = (r, T_n) .
\end{equation}
The series converges pointwise for functions which are continuous in $[-1, 1]$
up to a finite number of discontinuities in that interval. For the truncated
series
\begin{equation}
  r(x) \approx \sum_{n=0}^{N-1} c_n\,T_n(x),
\end{equation}
the error is spread smoothly over the interval $-1 \leq x \leq 1$. The
coefficients for the truncated series may be approximated by
\begin{equation}
  c_n = \frac{\pi}{N}\sum_{k=1}^{N} r(x_k)\,T_n(x_k),
  \qquad \text{where} \qquad
  x_k = \cos\!\left(\frac{\pi}{N}\left(k-\frac{1}{2}\right)\right) .
\end{equation}
Following [36, 60] we apply the method to computing the inverse square root of
$H^2$. It can be shown that for $|s| \leq 1$, $\|H\| \leq 8$ (in lattice
units). We denote the (in magnitude) smallest and the largest eigenvalues of
$H$ with $\alpha$ and $\beta$. Consequently $H^2$ has eigenvalues $\lambda$ in
an interval $[\alpha^2, \beta^2] \subset [0, 64]$. The general interval
$\lambda \in [\alpha^2, \beta^2]$ can be mapped into the generic domain
$x \in [-1, 1]$ by
\begin{equation}
  x = \frac{2\lambda - (\beta^2 + \alpha^2)}{\beta^2 - \alpha^2} .
\end{equation}
For the inverse square root function in question we then find
\begin{equation}
  r(x) = \left[\frac{1}{2}(\beta^2+\alpha^2) + \frac{1}{2}(\beta^2-\alpha^2)x\right]^{-1/2}
  = \frac{1}{\sqrt{\lambda(x)}},
\end{equation}
and with that function the coefficients are computed according to (7.89). The
approximation of the sign function is then obtained by multiplication with $H$:
\begin{align}
  \mathrm{sign}[H] &= \frac{H}{\sqrt{H^2}}
  = H\sum_{n=0}^{N-1} c_n\,T_n(X)
    + \order{\exp(-2N|\alpha/\beta|)}, \notag\\
  \text{with} \qquad
  X &\equiv \frac{2H^2 - (\beta^2+\alpha^2)\,\mathbb{1}}{\beta^2-\alpha^2} .
\end{align}
The error of the approximation series decreases
$\propto \exp(-2N|\alpha/\beta|)$, i.e., the number of terms necessary for a
requested accuracy grows proportional to $|\beta/\alpha|$. This ratio is just
the condition number of $H$: $|\lambda_{\max}/\lambda_{\min}|$. To improve the
condition number, the method is usually combined with a ``removal'' of small
eigenmodes of $H$, which we discuss below. The series may have several hundred
terms and therefore one should use error reducing techniques like the Clenshaw
recursion relation for summing it [61].

The polynomial approximation is technically simple and quite robust and may be
straightforwardly applied to the force calculations necessary when introducing
dynamical fermions (cf. Chap.~8). Theoretical arguments, however, suggest that
the polynomial approximation is less efficient than the rational approximation
methods and indeed it is more sensitive to the condition number. In a
comparison [58] for test examples the polynomial approximation typically
needed between 1.5 and 4 times more matrix--vector multiplications with the
operator $H$ than a rational approximation based on the Zolotarev method.

\textbf{Partial fractions \`a la Zolotarev:} In [57, 62] it was suggested to
approximate $1/\sqrt{x^2}$ by a ratio of polynomials. Meanwhile it has become
clear [58, 59, 63] that the best such approximation for the square root on an
interval $x^2 \in [1, \beta^2/\alpha^2]$ is the Zolotarev approximation [64].
In particular,
\begin{equation}
  \frac{1}{\sqrt{x^2}} \approx Q(x^2) = d\prod_{n=1}^{m}
  \frac{x^2 + c_{2n}}{x^2 + c_{2n-1}},
  \qquad
  c_n = \frac{\mathrm{sn}^2\left(nK(k')/(2m+1); k'\right)}
              {1 - \mathrm{sn}^2\left(nK(k')/(2m+1); k'\right)},
\end{equation}
with $k' = \sqrt{1-\alpha/\beta}$ and $K(k')$ is the so-called complete
elliptic integral and $\mathrm{sn}(u; k)$ denotes a Jacobi elliptic function
[65].

The normalization constant $d$ may be fixed by requiring $Q(1) = 1$. This
results in an approximation to the sign function that is bounded by 1 from
below (on the positive branch). Requesting instead that the approximation
fluctuates symmetrically around 1 corresponds to setting
\begin{equation}
  Q(1) + \frac{\beta^2}{\alpha^2}\,Q\!\left(\frac{\beta^2}{\alpha^2}\right) = 2 .
\end{equation}
The rational function can be decomposed into partial fractions such that
\begin{equation}
  Q(x^2) = d\,(x^2 + c_{2m})\sum_{n=1}^{m}\frac{b_n}{x^2 + c_{2n-1}},
\end{equation}
where the parameters $b_n$ may be computed from (7.93) giving
\begin{equation}
  b_n = \frac{\prod_{k=1}^{m-1}(c_{2k} - c_{2n-1})}
             {\prod_{k=1,\,k\neq n}^{m}(c_{2k-1} - c_{2n-1})} .
\end{equation}
In order to apply this to the matrix $H$ we have to scale it by dividing by
the absolute value of its smallest eigenvalue $|\alpha|$. The final formula for
the sign function, obtained by multiplication with $H$, thus reads
\begin{equation}
  \mathrm{sign}[H] \approx d\,\frac{H}{|\alpha|}\left[\left(\frac{H}{\alpha}\right)^2 + c_{2m}\right]
  \sum_{n=1}^{m} b_n\left[\left(\frac{H}{\alpha}\right)^2 + c_{2n-1}\right]^{-1} .
\end{equation}
Here $\alpha$ and $\beta$ denote the smallest and the largest eigenvalues of
$H$ and consequently the matrix $(H/\alpha)^2$ has its spectrum in
$[1, \beta^2/\alpha^2]$, as assumed for expansion (7.93).

Since in an application of (7.97) one always wants to compute the product of
the overlap matrix with some vector, one may compute the terms of (7.97)
simultaneously by a multi-shift conjugate gradient solver [66--68], as
discussed already in Chap.~6 in the context of computing the quark propagators
for several masses at once. The necessary overhead is small such that the
computational effort is essentially governed by the smallest $|c_{2n-1}|$.

Compared to other rational approximations, for a given accuracy, the Zolotarev
method drastically reduces the number of necessary terms in (7.97).

\textbf{Small eigenmode reduction:} In particular the polynomial approximation
is quite sensitive to the condition number of $H$. For $|s| < 1$ the
eigenvalues of $H$ are bounded from above, $|\lambda_{\max}| \leq 8$ (lattice
units), and therefore the true source of trouble is the smallest eigenvalue,
as discussed earlier. There is a method, however, to deal with cases with very
small $|\lambda_{\min}|$. We know from the spectral theorem (see
Appendix~A.5) that
\begin{align}
  f[H] &= \sum_i f(\lambda_i)\,v_i v_i^\dagger \notag\\
  &= \sum_{\lambda_i \geq \lambda_c} f(\lambda_i)\,v_i v_i^\dagger
    + \sum_{\lambda_i < \lambda_c} f(\lambda_i)\,v_i v_i^\dagger \notag\\
  &\equiv \sum_{\lambda_i < \lambda_c} f(\lambda_i)\,v_i v_i^\dagger
    + f\left[H^{\mathrm{(red)}}\right].
\end{align}
Here $v_i$ denotes the eigenvector for eigenvalue $\lambda_i$. We may
therefore split the problem into two parts. First we determine all eigenvalues
below some value $\lambda_c$ as well as their eigenvectors. Then we define a
reduced matrix by removing the small eigenvalue sector explicitly:
\begin{equation}
  H^{\mathrm{(red)}} \equiv H - \sum_{\lambda_i < \lambda_c}
  \lambda_i\,v_i v_i^\dagger .
\end{equation}
We compute $\mathrm{sign}\,H^{\mathrm{(red)}}$ for this reduced operator,
which has improved condition number, by some approximation scheme and invoke
(7.98) to reconstruct
\begin{equation}
  \mathrm{sign}[H] = \mathrm{sign}\left[H^{\mathrm{(red)}}\right]
  + \sum_{\lambda_i < \lambda_c}\mathrm{sign}(\lambda_i)\,v_i v_i^\dagger .
\end{equation}
The overhead due to the necessity to compute the low-eigenvalue sector is not
as disastrous as it may seem. This has to be done only once for a given gauge
configuration. The multiplication of the overlap Dirac operator (for that gauge
configuration) with a vector, however, has to be computed quite often, as it is
usually embedded in, e.g., a conjugate gradient solver used for calculating
the quark propagator. The effort of computing the, say, lowest 20 eigenvalues
is therefore in most situations small in comparison to the total cost.

\section*{References}

\begin{enumerate}[label={[\arabic*]}]
  \item J.~S.~Bell and R.~Jackiw: Nuovo Cimento 60A, 47 (1969).
  \item S.~L.~Adler: Phys.~Rev.~177, 2426 (1969).
  \item C.~Itzykson and J.-B.~Zuber: \emph{Quantum Field Theory} (McGraw-Hill,
        New York 1985).
  \item M.~E.~Peskin and D.~V.~Schroeder: \emph{An Introduction to Quantum
        Field Theory} (Addison-Wesley, Reading, Massachusetts 1995).
  \item H.~B.~Nielsen and M.~Ninomiya: Phys.~Lett.~B 105, 219 (1981).
  \item H.~B.~Nielsen and M.~Ninomiya: Nucl.~Phys.~B 185, 20 (1981).
  \item H.~B.~Nielsen and M.~Ninomiya: Nucl.~Phys.~B 193, 173 (1981).
  \item L.~H.~Karsten: Phys.~Lett.~B 192, 315 (1981).
  \item P.~H.~Ginsparg and K.~G.~Wilson: Phys.~Rev.~D 25, 2649 (1982).
  \item P.~Hasenfratz: Nucl.~Phys.~B (Proc.~Suppl.) 63A--C, 53 (1998).
  \item H.~Neuberger: Phys.~Lett.~B 427, 353 (1998).
  \item M.~L\"uscher: Phys.~Lett.~B 428, 342 (1998).
  \item S.~Chandrasekharan: Phys.~Rev.~D 60, 074503 (1999).
  \item M.~Atiyah and I.~M.~Singer: Ann.~Math.~93, 139 (1971).
  \item P.~Hasenfratz, V.~Laliena, and F.~Niedermayer: Phys.~Lett.~B 427, 125 (1998).
  \item A.~A.~Belavin, A.~M.~Polyakov, A.~S.~Schwartz, and Y.~S.~Tyupkin:
        Phys.~Lett.~B 59, 85 (1975).
  \item T.~Sch\"afer and E.~Shuryak: Rev.~Mod.~Phys.~70, 323 (1998).
  \item D.~Diakonov: Prog.~Part.~Nucl.~Phys.~51, 173 (2003).
  \item R.~Rajaraman: \emph{Solitons and Instantons} (North Holland Publishers,
        Amsterdam 1982).
  \item Y.~Kikukawa and A.~Yamada: Phys.~Lett.~B 448, 265 (1999).
  \item K.~Fujikawa: Nucl.~Phys.~B 546, 480 (1999).
  \item H.~Leutwyler and A.~Smilga: Phys.~Rev.~D 46, 5607 (1992).
  \item C.~Gattringer, R.~Hoffmann, and S.~Schaefer: Phys.~Lett.~B 535, 358 (2002).
  \item L.~Del Debbio and C.~Pica: JHEP 0402, 003 (2004).
  \item L.~Del Debbio, L.~Giusti, and C.~Pica: Phys.~Rev.~Lett.~94, 032003 (2005).
  \item E.~Witten: Nucl.~Phys.~B 156, 269 (1979).
  \item G.~Veneziano: Nucl.~Phys.~B 159, 213 (1979).
  \item G.~Veneziano: Phys.~Lett.~B 95, 90 (1980).
  \item L.~Giusti, G.~Rossi, M.~Testa, and G.~Veneziano: Nucl.~Phys.~B 628, 234 (2002).
  \item E.~Seiler: Phys.~Lett.~B 525, 355 (2002).
  \item L.~Giusti, G.~Rossi, and M.~Testa: Phys.~Lett.~B 587, 157 (2004).
  \item M.~L\"uscher: Phys.~Lett.~B 593, 296 (2004).
  \item M.~Aguado and E.~Seiler: Phys.~Rev.~D 72, 094502 (2005).
  \item J.~Osborn, D.~Toublan, and J.~Verbaarschot: Nucl.~Phys.~B 540, 317 (1999).
  \item P.~H.~Damgaard, J.~C.~Osborn, D.~Toublan, and J.~J.~Verbaarschot:
        Nucl.~Phys.~B 547, 305 (1999).
  \item P.~Hern\'andez, K.~Jansen, and L.~Lellouch: Phys.~Lett.~B 469, 198 (1999).
  \item P.~Hasenfratz et al.: Nucl.~Phys.~B 643, 280 (2002).
  \item T.~DeGrand: Phys.~Rev.~D 63, 034503 (2000).
  \item L.~Giusti, C.~Hoelbling, and C.~Rebbi: Phys.~Rev.~D 64, 114508 (2001).
  \item T.~Banks and A.~Casher: Nucl.~Phys.~B 169, 103 (1980).
  \item F.~Constantinescu: \emph{Distributions and Their Applications in
        Physics} (Pergamon Press, Oxford 1980).
\end{enumerate}
